\documentclass{article}%
\usepackage[T1]{fontenc}%
\usepackage[utf8]{inputenc}%
\usepackage{lmodern}%
\usepackage{textcomp}%
\usepackage{lastpage}%
\usepackage{geometry}%
\geometry{margin=1in,headheight=14pt,headsep=25pt}%
\usepackage{hyperref}%
\usepackage{enumitem}%
\usepackage{fancyhdr}%
\usepackage{xcolor}%
%

            \hypersetup{
                colorlinks=true,
                linkcolor=blue,
                filecolor=magenta,
                urlcolor=blue,
            }
            
            \pagestyle{fancy}
            \fancyhf{}
            \rhead{Generated on \today}
            \lhead{Course Summary}
            \cfoot{\thepage}
            
            % Configure itemize settings globally
            \setlist[itemize]{nosep}  % Reduce vertical spacing between items
            \setlist[itemize]{leftmargin=*}  % Align with left margin
        %
\title{Linear Programming Course Summary}%
\author{Generated by Scribe.AI}%
\date{\today}%
%
\begin{document}%
\normalsize%
\maketitle%
\section{Key Terms}%
\label{sec:KeyTerms}%
\subsection{Linear Programming Fundamentals}%
\label{subsec:LinearProgrammingFundamentals}%
\begin{itemize}[leftmargin=*]%
\item \href{https://www.math.purdue.edu/~yipn/421/2024-08-27-ExSimplex.pdf#page=1}{\textbf{Slack Variable}}: A variable added to a less-than-or-equal-to inequality constraint to transform it into an equality constraint. It represents the difference between the left-hand side and right-hand side of the inequality.%
\item \href{https://www.math.purdue.edu/~yipn/421/2024-08-27-ExSimplex.pdf#page=4}{\textbf{Feasible Region}}: The set of all points that satisfy all the constraints of a linear programming problem.%
\item \href{https://www.math.purdue.edu/~yipn/421/2024-08-27-ExSimplex.pdf#page=5}{\textbf{Vertex}}: A point where two or more constraint boundaries intersect in the feasible region of a linear programming problem. In the Simplex method, optimal solutions are found at vertices.%
\end{itemize}

%
\subsection{Simplex Method Algorithms}%
\label{subsec:SimplexMethodAlgorithms}%
\begin{itemize}[leftmargin=*]%
\item \href{https://www.math.purdue.edu/~yipn/421/2024-08-27-ExSimplex.pdf#page=6}{\textbf{Basic Variable}}: In the Simplex method, a variable that is expressed in terms of non-basic variables in a dictionary. In a basic feasible solution, basic variables are non-negative.%
\item \href{https://www.math.purdue.edu/~yipn/421/2024-08-27-ExSimplex.pdf#page=6}{\textbf{Non-basic Variable}}: In the Simplex method, a variable set to zero in a given iteration of the algorithm. The values of the basic variables are determined by the values of the non-basic variables.%
\item \href{https://www.math.purdue.edu/~yipn/421/2024-08-27-ExSimplex.pdf#page=9}{\textbf{Pivot Operation}}: The process of exchanging a basic and a non-basic variable in the Simplex method, leading to a new basic feasible solution.%
\item \href{https://www.math.purdue.edu/~yipn/421/2024-08-27-ExSimplex.pdf#page=26}{\textbf{Dictionary}}: A representation of a linear programming problem where basic variables are expressed as functions of non-basic variables. It is used in the Simplex method to iteratively find the optimal solution.%
\item \href{https://www.math.purdue.edu/~yipn/421/2024-08-27-ExSimplex.pdf#page=32}{\textbf{Auxiliary Problem}}: A modified linear programming problem introduced to find an initial feasible solution when the origin is not feasible in the original problem. It involves adding a new variable to shift the feasible region.%
\item \href{https://www.math.purdue.edu/~yipn/421/2024-09-05-SimplexEfficiency.pdf#page=4}{\textbf{Largest-Coefficient Rule}}: A pivot rule in the Simplex method that selects the entering variable with the largest coefficient in the objective function row of the dictionary.%
\end{itemize}

%
\section{Problem Types}%
\label{sec:ProblemTypes}%
\subsection{Simplex Method Analysis}%
\label{subsec:SimplexMethodAnalysis}%
\begin{itemize}[leftmargin=*]%
\item \href{https://www.math.purdue.edu/~yipn/421/2024-09-05-SimplexEfficiency.pdf#page=2}{\textbf{Worst-Case Analysis of the Simplex Method}}: The Klee-Minty example demonstrates that the Simplex method can require an exponential number of iterations in the worst case, depending on the pivot rule used.%
\item \href{https://www.math.purdue.edu/~yipn/421/2024-09-05-SimplexEfficiency.pdf#page=3}{\textbf{Impact of Pivot Rules on Simplex Method Efficiency}}: Different pivot rules (e.g., largest-coefficient, largest-increase) significantly affect the number of iterations required by the Simplex method, with some rules exhibiting exponential worst-case behavior while others perform better in practice.%
\item \href{https://www.math.purdue.edu/~yipn/421/2024-09-05-SimplexEfficiency.pdf#page=9}{\textbf{Empirical Analysis of Simplex Method Performance}}: Empirical studies show that the number of iterations in the Simplex method tends to grow polynomially with the problem size in practice, even though worst-case examples exist.%
\item \href{https://www.math.purdue.edu/~yipn/421/2024-09-05-SimplexEfficiency.pdf#page=10}{\textbf{Theoretical Bounds on Simplex Method Iterations}}: Theoretical analysis of randomized pivot rules provides bounds on the expected number of iterations, which are better than exponential but not polynomial.%
\item \href{https://www.math.purdue.edu/~yipn/421/2024-09-05-SimplexEfficiency.pdf#page=11}{\textbf{Smoothed Complexity of the Simplex Method}}: The smoothed complexity analysis suggests that small random perturbations to the input data can significantly improve the performance of the Simplex method, leading to a polynomial number of iterations with high probability.%
\end{itemize}

%
\subsection{Simplex Method Mechanics}%
\label{subsec:SimplexMethodMechanics}%
\begin{itemize}[leftmargin=*]%
\item \href{https://www.math.purdue.edu/~yipn/421/2024-08-27-ExSimplex.pdf#page=32}{\textbf{Finding an Initial Feasible Solution}}: The Simplex method requires an initial feasible solution. If the origin is not feasible, an auxiliary problem is introduced to find a feasible solution to the original problem.%
\item \href{https://www.math.purdue.edu/~yipn/421/2024-08-27-ExSimplex.pdf#page=30}{\textbf{Handling Infeasible Origins}}: The origin $(0, 0)$ may not satisfy all constraints in a linear programming problem. Techniques like introducing an auxiliary problem are used to find a feasible starting point.%
\item \href{https://www.math.purdue.edu/~yipn/421/2024-08-27-ExSimplex.pdf#page=17}{\textbf{Optimizing in Higher Dimensions}}: The Simplex method is extended to linear programming problems with more than two variables, requiring iterative pivoting operations to find the optimal solution.%
\item \href{https://www.math.purdue.edu/~yipn/421/2024-08-27-ExSimplex.pdf#page=26}{\textbf{Using Dictionaries to Represent Solutions}}: The Simplex method uses dictionaries to represent the current solution and to perform pivot operations efficiently. Each dictionary corresponds to a vertex of the feasible region.%
\end{itemize}

%
\section{Algorithm Solutions}%
\label{sec:AlgorithmSolutions}%
\subsection{Linear Programming Fundamentals}%
\label{subsec:LinearProgrammingFundamentals}%
\begin{itemize}[leftmargin=*]%
\item \href{https://www.math.purdue.edu/~yipn/421/2024-08-27-ExSimplex.pdf#page=6}{\textbf{Simplex Method}}: Iteratively move from one vertex of the feasible region to another by setting non-basic variables to zero and choosing a new set of (non)basic variables to improve the objective function value. This is equivalent to moving from vertex to vertex in the feasible region.%
\item \href{https://www.math.purdue.edu/~yipn/421/2024-08-27-ExSimplex.pdf#page=32}{\textbf{Auxiliary Problem}}: When the origin is not feasible, introduce a new variable $x_0$ and modify the constraints to create an auxiliary problem that maximizes $-x_0$. Solve this auxiliary problem to find a feasible solution for the original problem. If the optimal solution has $x_0 = 0$, then a feasible solution to the original problem is found.%
\end{itemize}

%
\subsection{Simplex Method Variants}%
\label{subsec:SimplexMethodVariants}%
\begin{itemize}[leftmargin=*]%
\item \href{https://www.math.purdue.edu/~yipn/421/2024-09-05-SimplexEfficiency.pdf#page=4}{\textbf{Largest-Coefficient Rule}}: Choose the variable with the largest coefficient in the objective function row to enter the basis.%
\item \href{https://www.math.purdue.edu/~yipn/421/2024-09-05-SimplexEfficiency.pdf#page=6}{\textbf{Largest-Increase Rule}}: Choose the variable whose entrance into the basis brings about the largest increase in the objective function.%
\item \href{https://www.math.purdue.edu/~yipn/421/2024-09-05-SimplexEfficiency.pdf#page=10}{\textbf{Bland's Rule}}: Choose the variable with the smallest index among those eligible to enter the basis.%
\item \href{https://www.math.purdue.edu/~yipn/421/2024-09-05-SimplexEfficiency.pdf#page=10}{\textbf{Randomized Pivot Rule}}: Randomly permute the indices of the variables; then use Bland's rule for entering and the lexicographic method for leaving variables.%
\end{itemize}

%
\subsection{Duality Theory}%
\label{subsec:DualityTheory}%
\begin{itemize}[leftmargin=*]%
\item \href{https://www.math.purdue.edu/~yipn/421/2024-09-12-Duality.pdf#page=2}{\textbf{Weak Duality}}: $c^T x \le b^T y$%
\item \href{https://www.math.purdue.edu/~yipn/421/2024-09-12-Duality.pdf#page=11}{\textbf{Strong Duality}}: If (P) has an optimal solution $x^$, then (D) has an optimal solution $y^$, and $c^T x^ = b^T y^$.%
\item \href{https://www.math.purdue.edu/~yipn/421/2024-09-12-Duality.pdf#page=21}{\textbf{Complementary Slackness}}: Primal feasible $x$ and dual feasible $y$ are optimal if and only if $x_j z_j = 0$ for all $j$ and $w_i y_i = 0$ for all $i$.%
\end{itemize}

%
\end{document}